\paragraph{目的限定と最小化}
分析目的に必要な列のみを用い、保存期間・アクセス権は社内規程に従う（詳細は \texttt{data\_dictionary.md}）。

\paragraph{説明可能性と人の関与}
セグメント別施策はモデル出力に基づくが、最終配信判断には担当者レビューを置く。自動化された決定のみで不利益を生じさせない運用を推奨。

\paragraph{差別・不利益リスク}
地域（\texttt{zip\_code} 等）などのプロキシ特性に依存する政策は、属性に基づく不当な差別に繋がらないよう設計・監査する。

\paragraph{オプトアウト・問い合わせ}
マーケティング配信の拒否・個人情報の開示等の請求手続きは、社内ポリシーおよび個人情報保護法（日本）の枠組みに従う。

\paragraph{法的助言ではない}
本稿・生成テンプレは\textbf{法的助言を構成しない}。実装・運用前に法務・DPO と整合を取ること。
